\documentclass[11pt]{article}

\usepackage[margin=1.1in]{geometry}
\usepackage{amsmath,amssymb,amsthm}
\usepackage{booktabs}
\usepackage{microtype}
\usepackage{xcolor}
\usepackage[colorlinks=true,linkcolor=blue!60!black,citecolor=blue!60!black,urlcolor=blue!60!black]{hyperref}

\newtheorem{theorem}{Theorem}[section]
\newtheorem{proposition}[theorem]{Proposition}
\newtheorem{corollary}[theorem]{Corollary}
\newtheorem{lemma}[theorem]{Lemma}
\theoremstyle{definition}
\newtheorem{definition}[theorem]{Definition}
\newtheorem{claim}[theorem]{Claim}
\theoremstyle{remark}
\newtheorem{remark}[theorem]{Remark}

\newcommand{\PT}{\mathcal{PT}}
\newcommand{\ii}{\mathrm{i}}
\newcommand{\ee}{\mathrm{e}}
\newcommand{\Tr}{\mathrm{Tr}}
\newcommand{\dd}{\mathrm{d}}
\newcommand{\norm}[1]{\left\lVert #1 \right\rVert}
\newcommand{\wt}{\widetilde{k}}

\title{Every exorcism has a price:\\
ghost removal as survival conditioning,\\
from Maldacena's boundary condition to the pole sheet}

\author{Jonathan Torkelson\\[2pt]
\small Embeddetech, Inc., Tulsa, Oklahoma, USA\\
\small\texttt{jtorkelson@embeddetech.com}}
\date{July 19, 2026}

\begin{document}
\maketitle

\begin{abstract}
\noindent\textbf{Draft --- not for circulation.}\medskip

\noindent
A higher-derivative (ghost) theory can be made ghost-free by a
choice: a boundary condition, a contour, a norm. The companion papers
of this series established a currency in which such choices can be
priced: a quantum evolution admits an ordinary
(``beable'') sub-stochastic process whenever its survival amplitude
stays bounded---divisor negativity then marking indivisibility rather
than forcing conditioning---and only when the amplitude grows without
bound does the dynamics survive solely as a \emph{survival-conditioned}
process---which, for the coupled ghost realizations priced here, is the
terminal-indivisibility class, with a strictly negative generator
floor (unboundedness alone forces conditioning but not terminality: a
decoupled growing mode stays divisible). This paper applies that currency to ghost \emph{removal}. The
thesis: every standard exorcism---boundary conditions that project the
ghost out, and propagator prescriptions that redefine its pole---is
survival conditioning in disguise, or buys its way out in another
trilemma currency; for the boundary-condition removals and the
Feynman route that price is computed here.
Three results. (i)~\emph{Maldacena's future boundary condition}
(conformal $\to$ Einstein gravity) is, mode by mode, a Gaussian effect
$\mathcal{M}$ (regularized $E_{\sigma}$) on the Ostrogradski velocity variable of a detuned
Pais--Uhlenbeck proxy whose detuning $\delta$ represents the de~Sitter
$\Delta=3/\Delta=2$ splitting. Its price sheet is closed form:
survival $p=[1+1/(2\sigma^{2}\delta)]^{-1/2}\xrightarrow{\sigma\to0}\sigma\sqrt{2\delta}$,
and Einstein-vacuum fidelity of the delivered state
$F=2\sqrt{\alpha\omega_{1}}/(\alpha+\omega_{1})\simeq2\sqrt{\omega_{2}\delta}$ (with
$\alpha=\omega_{1}\omega_{2}\delta$, in units $\omega_{1}=1$).
Both collapse at the conformal (Minkowski, $\delta\to0$) point---the
known flat-space obstruction, now priced; the convention-free content
is the survival collapse and the delivered state's absolute
$\delta^{-1/2}$ width divergence, the fidelity reading the latter
against a fixed Einstein-vacuum reference. (ii)~The
\emph{Hell--L\"ust--Zoupanos two-condition variant} $\mathcal{M}'$ fixes a
canonically conjugate pair ($\dot h$ and the conjugate combination
$\ddot h+\omega_{1}^{2}h$, $[q_{2},O_{2}]=\ii$), so $\mathcal{M}'$ is Heisenberg-limited: it
pays the same leading $\sigma\sqrt{2\delta}$ survival power plus an
$O(1)$ ordering/phase-space-cell surcharge, and what the surcharge
buys is not a cheaper exorcism but a more general one---the full
two-parameter Einstein solution with no in-state restriction. The
$\delta\to0$ collapse is not evaded. (iii)~\emph{The pole sheet is
readable in the divisor}: by Bochner's theorem a second-sheet
resonance (positive spectral measure) has $|A(t)|\le1$ and survives as
an \emph{unconditioned} sub-stochastic process (its divisor negativity
transient when the amplitude is eventually monotone decreasing), while a
first-sheet ghost pair (a discrete
complex-conjugate pole pair off the real axis) has unbounded $|A(t)|$
in states overlapping the growing leg, so conditioning
is \emph{mandatory}; for the dimer realization it is terminal, with the
generator floor of the first companion paper. Buoninfante's multi-particle masking law
$\sinh^{2}(\Gamma t/2)$ is exactly the time dependence of the
cross-population of the terminal-class (conditioning-mandatory) dimer
(with $\wt=\Gamma/2$), up
to the normalization $(\kappa/\wt)^{2}$, and his ghost-prescription
taxonomy (Feynman
ghost / anti-Feynman / fakeon) maps onto the trilemma currencies of
the series: sample space, causal order, silence. Unstable particles
survive unconditioned; ghosts survive only conditioned---and every
way of pretending otherwise pays, measurably, in one currency or
another.
\end{abstract}

\newpage
\begin{quote}
\small\noindent\textbf{Note on provenance.} This manuscript is an
experiment in AI-conducted research as much as a physics paper. The
project began when the author asked the large language model Claude
Fable~5 (Anthropic): \emph{``Is it not possible that the antilinear
PT/time-reversal symmetry element could maybe predict indivisible
non-Markovian processes?''} Everything that followed was performed by
the model: the research program, the theorems and proofs, the numerical
experiments, the literature search, and this text. The model was
repeatedly prompted to attack its own results adversarially and to
search for prior work. The author, a professional engineer rather than
a physicist, directed and reviewed the sessions but makes no claim to
comprehension of the findings, and is circulating the manuscript for
independent review by professional physicists. Accordingly, nothing
here should be taken on trust: every result is stated with the proof
it actually has---closed-form derivations where they exist, and
numerically established claims labeled as such---and every numerical
claim is reproducible from the public scripts (see Code availability);
the originating conversation, project briefing, and full research
notebook are preserved in the program's hub repository,
\url{https://github.com/embeddetech/ghosts-program}.
\end{quote}

\newpage
\tableofcontents

%=====================================================================
\section{Introduction}
\label{sec:intro}

Ghosts are removed, in the literature, in two broad ways. The first
is by \emph{boundary condition}: one imposes conditions at a future
(or asymptotic) boundary that the ghost modes fail, so that only the
healthy sector survives. The paradigm is Maldacena's derivation of
Einstein gravity from conformal gravity \cite{Maldacena2011}: the
fourth-order graviton equation has four solutions per mode, a
positive-frequency condition in the far past halves them, and a
Neumann condition $\dot g=0$ at the future boundary halves them
again, leaving exactly the Einstein graviton; at tree level the
conformal-gravity wavefunctional evaluated on Neumann data reproduces
the (renormalized) Bunch--Davies wavefunctional of Einstein gravity.
A recent variant by Hell, L\"ust, and Zoupanos \cite{HLZ} keeps the
Neumann condition but adds a second future condition on $\ddot h$ and
drops the in-state restriction, aiming at the \emph{general} Einstein
solution. The second removal is by \emph{prescription}: one keeps the
ghost in the propagator but redefines how its pole is handled---the
Feynman ghost, the anti-Feynman (advanced-arrow) ghost, and Anselmi's
fakeon are the three canonical choices, recently organized into a
sharp taxonomy by Buoninfante \cite{BuoninfanteP1,BuoninfanteP2,%
BuoninfanteP3} and disputed by Anselmi, Briscese, Calcagni, and
Modesto \cite{Anselmi2025}, with the Donoghue--Menezes ``Merlin
mode'' treatment \cite{DonoghueMenezes} occupying a related corner.

The companion papers of this series supply a currency in which such
removals can be priced rather than merely asserted. The foundational
paper \cite{indivisible} shows that a non-unitary or higher-derivative
quantum evolution admits an ordinary (``beable'') sub-stochastic
process---definite configurations, nonnegative transition
probabilities---whenever its survival amplitude stays bounded; the
two-time \emph{divisor}
$M(t_{2},t_{1})=\Gamma(t_{2})\Gamma(t_{1})^{-1}$ of its transition
matrices staying entrywise nonnegative on $[t_{1},t_{2}]$ marks the
process \emph{divisible} there, and divisor negativity marks
indivisibility, not non-existence. When the negativity is only
transient, the process is indivisible but ordinary. Only when the
amplitude grows without bound---so that no unconditioned sub-stochastic
process exists---does the dynamics survive solely as a
\emph{survival-conditioned} process, which for the coupled ghost
realizations treated here is the terminal-indivisibility class
(permanent divisor negativity past a critical time, the conditioned
generator pinned at a strictly negative floor), in which every
probability is conditioned on an ensemble that is itself being consumed.
(Unboundedness by itself forces the conditioning but not the
terminality---a decoupled mode that merely grows stays divisible; the
terminal floor is a feature of the coupling.) The
trilemma companion \cite{choosing} then prices the field-theory
options: a coupled-ghost theory keeps at most two of \{coupled
ghosts, empirical Lorentz invariance, definite configurations\} at
tree level, and \cite{instantaneous} makes the Lorentz cost concrete
through the boost-orbit measure of ghost vacuum decay.

The Ostrogradski ghost of a higher-derivative theory is, on this
view, not a calculational error to be deleted but a structural
feature that must be paid for. This paper's thesis is that ghost
removal itself belongs on that price sheet:
\begin{quote}
\emph{Every standard exorcism is survival conditioning in disguise,
or buys its way out in another trilemma currency; for the
boundary-condition removals and the Feynman route that price is
computed here.}
\end{quote}
For boundary-condition removals the statement is nearly literal: a
future boundary condition that the prepared state does not already
satisfy is a quantum \emph{effect}---a post-selection---and the
process that results is survival-conditioned by construction. What
has been missing is the price. For prescription removals the
statement is structural: the divisor reads the pole sheet, and the
first-sheet ghost pole forces conditioning while the second-sheet
resonance does not; the prescriptions that avoid conditioning pay
instead in the other trilemma currencies.

Concretely, we establish three groups of results.

\paragraph{(i) Maldacena's boundary condition, priced
(Section~\ref{sec:maldacena}).}
We translate the per-mode fourth-order graviton system into a detuned
Pais--Uhlenbeck (PU) oscillator with frequencies
$\omega_{1}$ and $\omega_{2}=\omega_{1}-\delta$, the detuning
$\delta$ standing for the de~Sitter splitting of the
$\Delta=3$ graviton from its $\Delta=2$ negative-norm partner;
Minkowski is the conformal point $\delta\to0$, where the pair
degenerates into a Jordan block. In Ostrogradski variables
$q_{1}=x$, $q_{2}=\dot x$ the Neumann condition is a condition on a
\emph{configuration} variable, so its regularized form is a bona fide
multiplication effect
$E_{\sigma}=\exp(-q_{2}^{2}/2\sigma^{2})$, $0\le E_{\sigma}\le1$;
Maldacena's positive-frequency condition is the state preparation and
is deliberately \emph{not} counted in the price. On the
Dirac-normalizable vacuum of the declared quantization the price
sheet is closed form:
\[
p(\sigma,\delta)=\Bigl[1+\tfrac{1}{2\sigma^{2}\delta}\Bigr]^{-1/2}
\;\xrightarrow{\ \sigma\to0\ }\;\sigma\sqrt{2\delta},
\qquad
F=\frac{2\sqrt{\alpha\,\omega_{1}}}{\alpha+\omega_{1}}
\simeq2\sqrt{\omega_{2}\delta},
\quad \alpha=\omega_{1}\omega_{2}\delta \quad(\omega_{1}=1),
\]
where $p$ is the survival probability of the effect and $F$ the
fidelity of the delivered (conditioned) graviton state to the
Einstein vacuum (against a fixed unit-mass reference; the
convention-free content is the survival collapse and the delivered
state's absolute width divergence). Both vanish as $\delta\to0$: the exorcism works in
de~Sitter, where $\delta=O(1)$, and its price diverges exactly at the
Minkowski point where Hell--L\"ust--Zoupanos found the flat-space
version fails. The known qualitative obstruction becomes a
quantitative collapse of the price sheet.

\paragraph{(ii) The Hell--L\"ust--Zoupanos variant is
Heisenberg-limited (Section~\ref{sec:hlz}).}
HLZ impose two future conditions---Neumann $\dot h|_{0}=0$
(their Eq.~(30)) and a curvature-matching condition on $\ddot h$
(their Eq.~(59) in covariant form; per Fourier mode their Eq.~(64),
$\ddot h_{k}|_{0}=k^{2}h_{k}|_{0}$)---and no in-state restriction,
delivering the general two-parameter Einstein-$\Lambda$ tensor mode
(their Eq.~(67)). In the Ostrogradski frame the second condition
constrains $O_{2}=p_{2}+\omega_{1}^{2}q_{1}$, and $p_{2}$ is the
momentum \emph{conjugate} to the $q_{2}$ that Neumann fixes:
$[q_{2},O_{2}]=\ii$. The two-condition effect $\mathcal{M}'$ is therefore a
joint position--momentum (Heisenberg-limited) effect. Its
\emph{Weyl-ordered} survival factorizes at the Gaussian level
($\mathrm{cov}(q_{2},O_{2})=0$ exactly), the second condition is
asymptotically free
($\mathrm{Var}(O_{2})=\tfrac{\delta}{2}\,
\tfrac{\omega_{1}+\omega_{2}}{\omega_{2}}\to0$; the ghost-killing sign
is also the finite-variance one, the opposite sign giving
$\mathrm{Var}(O_{2})\to\infty$, Prop.~\ref{prop:O2}), and both natural
operator orderings then give the same leading power
$p_{\mathcal{M}'}\sim\sigma\sqrt{2\delta}$, differing only by an $O(1)$
phase-space-cell factor. The surcharge buys generality, not
cheapness; the conformal collapse is not evaded.

\paragraph{(iii) The pole sheet is readable in the divisor
(Sections~\ref{sec:sheet}--\ref{sec:taxonomy}).}
Buoninfante's trilogy \cite{BuoninfanteP1,BuoninfanteP2,BuoninfanteP3}
isolates the sharp spectral difference between an ordinary unstable
particle and a ghost with the same complex-conjugate pole pair
$m^{2}\pm\ii m\Gamma$: for the unstable particle the pair sits on the
\emph{second} sheet and the K\"all\'en--Lehmann measure is a positive
density with sum rule $1=C$ (all norm transfers to the continuum);
for the ghost the pair sits on the \emph{first} sheet and the
spectral decomposition acquires discrete \emph{negative-residue}
complex-pole terms with sum rule $Z+Z^{*}=1+C$ (the one-particle
component persists). We show this dichotomy is exactly the divisor
dichotomy of \cite{indivisible}: a positive spectral measure makes
the survival amplitude positive-definite in Bochner's sense, so
$|A(t)|\le1$ and an \emph{unconditioned} sub-stochastic survival
process exists, its divisor negativity transient when the amplitude is
eventually monotone decreasing; the
first-sheet pair makes $|A(t)|$ grow without bound, so no
unconditioned process exists and the survival-conditioned dilation is
mandatory---and, for the dimer realization, it lands in the terminal
class, with the critical time and generator floor of \cite{indivisible}
in closed form.
Buoninfante's ``anti-instability'' masking law, the multi-particle
overlap growing as $\sinh^{2}(\Gamma t/2)$, is shown to be
\emph{exactly} the time dependence of the first companion paper's
terminal-class (conditioning-mandatory) dimer cross-population---the
pre-conditioning growth that conditioning divides out---%
$|U_{12}(t)|^{2}
=(\kappa/\wt)^{2}\sinh^{2}(\wt t)$ with $\wt=\Gamma/2$, up to the
normalization $(\kappa/\wt)^{2}$: two
bookkeepings of the same pole pair, one in the closed indefinite-metric
currency, one in the beable currency. Finally the prescription
taxonomy maps onto the trilemma: the Feynman ghost pays sample space
(conditioning), the anti-Feynman ghost pays causal order (the
Donoghue--Menezes corner---a connection this paper draws), and the
fakeon pays silence (the state is excised). The live Anselmi--Buoninfante dispute over which
prescription is ``viable'' is, in this reading, a currency dispute,
and the trilemma prices both sides.

Section~\ref{sec:priorart} scopes the claims against the final-state
post-selection literature \cite{HorowitzMaldacena,Anastopoulos2024};
Section~\ref{sec:limitations} lists the limitations, all of which are
substantive; Section~\ref{sec:conclusion} concludes. Throughout,
statements proved in closed form are labeled propositions with
proofs; statements established numerically with analytic structure
are labeled as such in the proof text. All numbers are reproducible
from three short scripts (Code availability).

%=====================================================================
\section{The currency: divisors, conditioning, and the detuned
Pais--Uhlenbeck mode}
\label{sec:currency}

\subsection{Divisors and the terminal class}
\label{sec:divisor}

We recall the minimal apparatus of \cite{indivisible}. Let a
finite-configuration dynamics be presented by transition matrices
$\Gamma(t)$ from time $0$ to time $t$ (columns summing to at most
one: a \emph{sub-stochastic} process, the deficit being read as
decay or, in the conditioned case, as failure to survive
post-selection). The process is an ordinary (beable) process on the
interval $[t_{1},t_{2}]$ iff the \emph{divisor}
\begin{equation}
M(t_{2},t_{1})\;=\;\Gamma(t_{2})\,\Gamma(t_{1})^{-1}
\label{eq:divisor}
\end{equation}
is entrywise nonnegative: $M$ is then itself a transition matrix and
the process can be continued through $t_{1}$ without memory of how it
got there. (For a strictly sub-stochastic $\Gamma$ one additionally
needs $M$ column-sub-stochastic; entrywise nonnegativity is the
operative criterion, the normalization tracked separately, and the full
statement is that of \cite{indivisible}.) Transient divisor negativity (indivisibility) is
compatible with an ordinary process on coarser time grids; the
pathological class identified in \cite{indivisible} is
\emph{terminal} indivisibility: negativity that sets in at a finite
$t_{c}$ and never leaves, with the late-time conditioned generator
pinned at a strictly negative off-diagonal floor. A unitary (or
non-unitary) mode dynamics $U(t)$ that admits no sub-stochastic
$\Gamma$ at all---because $\norm{U(t)}>1$---survives in the beable
currency only through the survival-conditioned dilation: normalize by
the largest singular value,
\begin{equation}
\Gamma_{\mathrm{c}}(t)\;=\;
\Bigl|\,U(t)/\sigma_{\max}\!\bigl(U(t)\bigr)\Bigr|^{2}
\quad(\text{entrywise modulus squared}),
\label{eq:conditioned}
\end{equation}
which is the transition matrix of the process conditioned on
surviving the growth. For the broken-$\PT$ dimer
\begin{equation}
H=\begin{pmatrix}\ii a&\kappa\\ \kappa&-\ii a\end{pmatrix},
\qquad a>\kappa,\qquad \wt=\sqrt{a^{2}-\kappa^{2}},
\label{eq:dimer}
\end{equation}
the first companion paper derives the closed forms
\begin{equation}
t_{c}=\frac{1}{\wt}\,
\operatorname{arcsinh}\!\Bigl(\frac{\wt}{\kappa\sqrt{2}}\Bigr),
\qquad
\text{generator floor}\;=\;-(a+\wt),
\label{eq:tcfloor}
\end{equation}
both of which are re-verified numerically in
Section~\ref{sec:sheet}. Because $H^{2}=-\wt^{2}\mathbf{1}$
(the $\sigma_{z}$ and $\sigma_{x}$ parts anticommute),
\begin{equation}
U(t)=\ee^{-\ii Ht}
=\cosh(\wt t)\,\mathbf{1}-\ii\,\frac{\sinh(\wt t)}{\wt}\,H,
\qquad
|U_{12}(t)|^{2}
=\frac{\kappa^{2}}{\wt^{2}}\,\sinh^{2}(\wt t),
\label{eq:dimerU}
\end{equation}
an exact identity we will need in Section~\ref{sec:masking}.

\subsection{The detuned Pais--Uhlenbeck proxy}
\label{sec:pu}

Per mode, the linearized fourth-order graviton dynamics of conformal
gravity is a Pais--Uhlenbeck oscillator \cite{PaisUhlenbeck1950}. We
work with the detuned PU Lagrangian
\begin{equation}
L=\tfrac12\bigl[\ddot x^{2}
-(\omega_{1}^{2}+\omega_{2}^{2})\,\dot x^{2}
+\omega_{1}^{2}\omega_{2}^{2}\,x^{2}\bigr],
\qquad
\omega_{2}=\omega_{1}-\delta,\qquad 0<\delta<\omega_{1},
\label{eq:PUlag}
\end{equation}
with solutions $\ee^{\pm\ii\omega_{1}t},\ee^{\pm\ii\omega_{2}t}$. We
set $\omega_{1}=1$ throughout the numerics. The detuning $\delta$ is
our proxy dial for the de~Sitter splitting: in
Maldacena's analysis \cite{Maldacena2011} the surviving graviton is
the $\Delta=3$ mode and its removed partner is the negative-norm
$\Delta=2$ mode, split by the Hubble scale; the Minkowski limit
$H_{\Lambda}\to0$ merges the pair into the degenerate ($t$-growing,
Jordan-block) conformal mode, i.e.\ $\delta\to0$. The proxy is a
declared, mode-level translation---we do not derive it from the
conformal-gravity path integral, and Section~\ref{sec:limitations}
returns to this point.

Ostrogradski variables: $q_{1}=x$, $q_{2}=\dot x$, with conjugate
momenta $p_{1},p_{2}$ ($p_{2}=\ddot x$ on shell), and
\begin{equation}
H=p_{1}q_{2}+\tfrac12 p_{2}^{2}
+\tfrac12(\omega_{1}^{2}+\omega_{2}^{2})\,q_{2}^{2}
-\tfrac12\,\omega_{1}^{2}\omega_{2}^{2}\,q_{1}^{2}.
\label{eq:PUham}
\end{equation}
Both $q_{1}$ and $q_{2}$ are \emph{configuration} variables: the
beable sample space of this mode is the two-slice memory
$(x,\dot x)$, exactly as in the companion papers. $H$ is Hermitian on
$L^{2}(\mathbb{R}^{2})$ and unbounded below (Ostrogradski); its
normalizable spectrum in this quantization is
$E_{n_{1}n_{2}}=\omega_{1}(n_{1}+\tfrac12)
-\omega_{2}(n_{2}+\tfrac12)$.

\begin{proposition}[The Dirac-normalizable vacuum and its
exceptional-point degeneration]
\label{prop:vacuum}
The Gaussian ansatz
$\Psi_{0}(q_{1},q_{2})
=\exp\bigl[-\tfrac12(\alpha q_{1}^{2}+bq_{2}^{2})-\ii c\,q_{1}q_{2}\bigr]$
solves $H\Psi_{0}=E_{0}\Psi_{0}$ exactly, with
\begin{equation}
c=\omega_{1}\omega_{2},\qquad
b=\delta,\qquad
\alpha=\omega_{1}\omega_{2}\,\delta,\qquad
E_{0}=\tfrac{\delta}{2}.
\label{eq:vacuum}
\end{equation}
Hence $|\Psi_{0}|^{2}
=\exp(-\omega_{1}\omega_{2}\delta\,q_{1}^{2}-\delta\,q_{2}^{2})$: the
state is normalizable for every $\delta>0$, and its widths in
\emph{both} configuration variables diverge as $\delta^{-1/2}$ at the
conformal point. The root $c=-\omega_{1}\omega_{2}$,
$b=\omega_{1}+\omega_{2}$ has $\alpha=-\omega_{1}\omega_{2}
(\omega_{1}+\omega_{2})<0$ and is normalizable only under contour
rotation (the Bender--Mannheim quantization
\cite{BenderMannheim2008}); it is not used here.
\end{proposition}

\begin{proof}
Substituting the ansatz into \eqref{eq:PUham} with
$p_{1}=-\ii\partial_{q_{1}}$, $p_{2}=-\ii\partial_{q_{2}}$ and
collecting monomials gives three coefficient equations and the
eigenvalue:
\[
q_{1}^{2}:\ \ \tfrac{c^{2}}{2}
=\tfrac{\omega_{1}^{2}\omega_{2}^{2}}{2};
\qquad
q_{2}^{2}:\ \ \tfrac{b^{2}}{2}+c
=\tfrac{\omega_{1}^{2}+\omega_{2}^{2}}{2};
\qquad
q_{1}q_{2}:\ \ \alpha=bc;
\qquad
E_{0}=\tfrac{b}{2}.
\]
The first fixes $c=\pm\omega_{1}\omega_{2}$; with
$c=+\omega_{1}\omega_{2}$ the second gives
$b^{2}=(\omega_{1}-\omega_{2})^{2}$, and the normalizable root is
$b=\delta$, whence $\alpha=\omega_{1}\omega_{2}\delta$ and
$E_{0}=\delta/2$. With $c=-\omega_{1}\omega_{2}$ one finds
$b=\omega_{1}+\omega_{2}$, $\alpha<0$. The three independent coefficient
equations were
also confirmed to machine precision, and the eigenvalue on a
$1201^{2}$ grid at $\delta=0.5$ to $E_{0}=0.249998$ vs.\ $0.25$
(the script noted in Code availability).
\end{proof}

The $\delta^{-1/2}$ width divergence is the exceptional-point
degeneration of the equal-frequency limit made quantitative at the
level of the state: as the two frequencies merge, the normalizable
vacuum flattens in both beables at the same rate. Everything priced
below inherits this degeneration.

\paragraph{Declared quantization (and why it must be declared).}
The PU mode at unequal frequencies admits inequivalent quantizations:
the Dirac-normalizable realization above (positive norm, energy
unbounded below---the ``Hermitian ghost'' realization, which is the
one carrying a beable sample space, cf.\ \cite{choosing}), and the
contour-rotated Bender--Mannheim realization (bounded energy,
indefinite or redefined norm). The price computed in this paper is a
statement about the first. This is the honest order of operations:
one must choose the quantization \emph{before} pricing the effect,
and the answer may depend on the choice. We state it once and do not
hide it.

%=====================================================================
\section{Maldacena's exorcism, priced}
\label{sec:maldacena}

\subsection{What is the effect and what is the state preparation}
\label{sec:effect}

Maldacena's argument \cite{Maldacena2011} has two moving parts per
graviton mode. First, positive frequency in the far past: this
selects the Bunch--Davies-type in-state. Second, the Neumann
condition $\dot g=0$ at the future boundary: this kills the
$\Delta=2$ partner and leaves the Einstein graviton, with
$\Psi_{\mathrm{Conformal}}[g,\dot g=0]
\propto\Psi_{\mathrm{BD\text{-}ren}}[g]$ at tree level. In
post-selection language the two parts are \emph{not} on the same
footing, and pricing them as if they were would double-count:
\begin{itemize}
\item the positive-frequency condition is \emph{state preparation}
(it fixes $\rho_{\mathrm{in}}$, here proxied by
$\Psi_{0}$ of Proposition~\ref{prop:vacuum});
\item the Neumann condition is the \emph{effect}: a condition imposed
at a later boundary that the prepared state does not already satisfy.
\end{itemize}
Only the second carries a survival price, and only the second is what
we price.

The sharp Neumann effect ``$q_{2}=0$'' is improper (a
$\delta$-normalized projection, with survival probability strictly
zero against any absolutely continuous distribution). We therefore
price the regularized Gaussian family
\begin{equation}
E_{\sigma}=\exp\!\Bigl(-\frac{q_{2}^{2}}{2\sigma^{2}}\Bigr),
\qquad 0\le E_{\sigma}\le1,
\qquad K_{\sigma}=E_{\sigma}^{1/2},
\label{eq:Esigma}
\end{equation}
a bona fide multiplication effect on the configuration variable
$q_{2}=\dot x$ (this is the point of the Ostrogradski frame: the
velocity is a beable, so the Neumann condition is a position-type
measurement, not a momentum smearing), and control the
$\sigma\to0$ limit explicitly. The determinant of $E_{\sigma}$ on the
relevant Gaussian sector and the survival probability vanish
together, so the regularization introduces no spurious hierarchy;
all statements below are made at fixed $\sigma$ and then in the
stated $\sigma\to0$ asymptotics.

\subsection{The survival price}

\begin{proposition}[Survival]
\label{prop:survival}
On the vacuum of Proposition~\ref{prop:vacuum},
\begin{equation}
p(\sigma,\delta)
=\frac{\langle\Psi_{0}|E_{\sigma}|\Psi_{0}\rangle}
{\langle\Psi_{0}|\Psi_{0}\rangle}
=\Bigl[1+\frac{1}{2\sigma^{2}\delta}\Bigr]^{-1/2}
\;\xrightarrow{\ \sigma\to0\ }\;
\sigma\sqrt{2\delta}\,.
\label{eq:price}
\end{equation}
\end{proposition}

\begin{proof}
The $q_{2}$-marginal of $|\Psi_{0}|^{2}$ is
$\propto\ee^{-\delta q_{2}^{2}}$ (the $q_{1}$ integral factors), so
\[
p=\frac{\int\ee^{-\delta q_{2}^{2}}
\ee^{-q_{2}^{2}/2\sigma^{2}}\,\dd q_{2}}
{\int\ee^{-\delta q_{2}^{2}}\,\dd q_{2}}
=\sqrt{\frac{\delta}{\delta+1/2\sigma^{2}}}
=\Bigl[1+\tfrac{1}{2\sigma^{2}\delta}\Bigr]^{-1/2}.
\]
The grid evaluation agrees with this closed form to six digits at all
tested $(\sigma,\delta)$.
\end{proof}

The physics is in the $\delta$-dependence, and it is transparent: the
vacuum's spread in the very variable the exorcism constrains grows as
$\delta^{-1/2}$ (Proposition~\ref{prop:vacuum}), so the fraction of
the ensemble that survives the constraint shrinks as
$\sqrt{\delta}$. In de~Sitter, where the $\Delta=3/\Delta=2$
splitting is order one, the price is order $\sigma$---the ordinary
cost of any finite-resolution post-selection. As the cosmological
constant is switched off the exorcism becomes arbitrarily expensive.

\subsection{The delivered state}

Removing the ghost is only half of the advertisement; the other half
is that what survives is \emph{Einstein} gravity. The proxy for that
claim is the fidelity of the conditioned state to the Einstein-mode
vacuum $\phi(q_{1})=(\omega_{1}/\pi)^{1/4}
\ee^{-\omega_{1}q_{1}^{2}/2}$.

\begin{proposition}[Delivery]
\label{prop:delivery}
As $\sigma\to0$ the conditioned reduced state on $q_{1}$ is the pure
Gaussian $\psi_{c}(q_{1})\propto\Psi_{0}(q_{1},0)
=\ee^{-\alpha q_{1}^{2}/2}$ with $\alpha=\omega_{1}\omega_{2}\delta$, and its
fidelity to the Einstein vacuum is
\begin{equation}
F=\bigl|\langle\phi|\psi_{c}\rangle\bigr|^{2}
=\frac{2\sqrt{\alpha\,\omega_{1}}}{\alpha+\omega_{1}}
\;=\;\frac{2\sqrt{\omega_{2}\delta}}{1+\omega_{2}\delta}
\Big|_{\omega_{1}=1}
\;\simeq\;2\sqrt{\omega_{2}\delta}
\quad(\delta\to0).
\label{eq:fidelity}
\end{equation}
\end{proposition}

\begin{proof}
As $\sigma\to0$ the Kraus operator $K_{\sigma}$ concentrates the
$q_{2}$ dependence on $q_{2}=0$; since
$\Psi_{0}(q_{1},q_{2})$ is Gaussian, the conditioned $q_{1}$ slice is
$\ee^{-\alpha q_{1}^{2}/2}$ (the phase $\ee^{-\ii cq_{1}q_{2}}$ drops on the
slice). The overlap of two centered Gaussians of widths
$\alpha^{-1/2},\omega_{1}^{-1/2}$ gives \eqref{eq:fidelity} by a one-line
integral. Numerically, the log--log slope of $F(\delta)$ over
$\delta\in[8\times10^{-4},2\times10^{-2}]$ is $0.491$
(predicted $1/2$), and the finite-$\sigma$ \emph{mixed} conditioned
state at $(\delta,\sigma)=(0.1,0.05)$ gives $F=0.5500$ against the
$\sigma\to0$ closed form $0.5505$: the pure-slice limit is reached
smoothly.
\end{proof}

The delivered graviton is a Gaussian of width $\alpha^{-1/2}
\sim(\omega_{2}\delta)^{-1/2}$, not the Einstein-vacuum width
$\omega_{1}^{-1/2}$: the exorcised state remembers the
exceptional-point flattening of the parent vacuum. At $\delta=O(1)$
(de~Sitter) the fidelity is order one ($F=0.80$ at $\delta=0.5$); at
the conformal point it collapses like $\sqrt{\delta}$. This fidelity is
taken against a \emph{fixed} external Einstein vacuum $\phi$ of
$\delta$-independent width $\omega_{1}^{-1/2}$ (the healthy graviton at
unit mass)---a declared convention. The proxy's \emph{own} healthy
normal mode co-degenerates with the delivered state
($\mathrm{Var}(y_{1})=1/(2\omega_{1}\delta(\omega_{1}+\omega_{2}))\to\infty$
as $\delta\to0$), so measured against that co-scaling reference the
fidelity would stay $O(1)$ (tending to $2\sqrt{2}/3$). The
$\sqrt{\delta}$ collapse is therefore the delivered state's
\emph{absolute} width divergence $\alpha^{-1/2}\sim\delta^{-1/2}$ read
against a fixed ruler, not a reference-independent fact; the
convention-free content is (a)~the survival collapse
$p\sim\sigma\sqrt{2\delta}$ and (b)~that absolute width divergence. (At
$\delta=0$ neither state is normalizable, so $F$ is a ratio of two
diverging widths there.)

\subsection{The price sheet and the conformal collapse}
\label{sec:pricesheet}

\begin{table}[ht]
\centering
\begin{tabular}{ccc}
\toprule
$\delta$ & $p$ (survival, $\sigma=0.05$) & $F$ (Einstein fidelity) \\
\midrule
$0.500$ & $0.04994$ & $0.8000$ \\
$0.100$ & $0.02236$ & $0.5505$ \\
$0.020$ & $0.01000$ & $0.2746$ \\
$0.004$ & $0.00447$ & $0.1257$ \\
\bottomrule
\end{tabular}
\caption{The price sheet of Maldacena's exorcism on the detuned-PU
proxy: survival probability of the regularized Neumann effect
(Proposition~\ref{prop:survival}) and Einstein-vacuum fidelity of the
delivered state (Proposition~\ref{prop:delivery}), both closed form,
both grid-verified. Both columns collapse at the conformal
(Minkowski) point $\delta\to0$; the survival collapse and the delivered
state's \emph{absolute} width divergence are convention-free, while the
fidelity column reads that divergence against a \emph{fixed}
Einstein-vacuum width (text after Proposition~\ref{prop:delivery}).}
\label{tab:pricesheet}
\end{table}

\begin{corollary}[Conformal collapse]
\label{cor:collapse}
$p\,F\simeq2\sqrt{2}\,\sigma\,\omega_{2}^{1/2}\,\delta$ as
$\delta\to0$ at fixed $\sigma$: the joint figure of merit of the
exorcism---how much ensemble survives times how Einstein-like the
survivor is---vanishes \emph{linearly} in the de~Sitter splitting. (The
$F$ factor inherits the fixed-reference convention of
Proposition~\ref{prop:delivery}; the survival factor $p$ is
convention-free.)
\end{corollary}

This is the quantitative form of a failure that Hell, L\"ust, and
Zoupanos located qualitatively \cite{HLZ}. Their flat-space analysis
finds that the Neumann-selected survivor in Minkowski is \emph{not}
an Einstein solution---a $t\,\ee^{-\ii kt}$ growing piece survives
inside the selected state---while in de~Sitter the survivor is
Einstein but carries only half the Einstein data. Our diagnosis of
the root cause: flat space is the degenerate PU point. At
$\delta=0$ the four-dimensional solution space does not split
invariantly into graviton~$\otimes$~partner (the mode matrix is a
Jordan block), so writing the removal as
$\mathcal{M}=\mathbf{1}\otimes E$ is singular exactly at $H_{\Lambda}=0$;
de~Sitter works because the $\Delta=3/\Delta=2$ splitting resolves
the degeneracy. The prediction that follows---the price should
degenerate as $H_{\Lambda}\to0$---is what
Table~\ref{tab:pricesheet} confirms: the divisor currency
\emph{detects} the flat-space obstruction as a price collapse, without
being told about Jordan blocks. For the de~Sitter over-selection
(half the Einstein data), the post-selection reading is actually
charitable: an effect is supposed to pin a state; the ``missing''
half of the classical data is not a bug of the exorcism-as-conditioning
but of reading it as a classical solution map. What one \emph{does}
pay to recover the general Einstein solution is the subject of the
next section.

%=====================================================================
\section{The Hell--L\"ust--Zoupanos variant: a Heisenberg-limited
exorcism}
\label{sec:hlz}

\subsection{Two future conditions, one conjugate pair}

HLZ \cite{HLZ} replace Maldacena's (state-prep $+$ one condition)
scheme by two future conditions and no in-state restriction: the
Neumann condition $\dot h|_{0}=0$ (their Eq.~(30)) and the
curvature-matching condition on the second derivative (their
Eq.~(59) in covariant form), which for a de~Sitter tensor mode reads
$\ddot h_{k}|_{0}=k^{2}h_{k}|_{0}$ (their Fourier-space Eq.~(64));
together these fix the secular (ghost) coefficients in terms of the
oscillatory ones and deliver the \emph{general} two-parameter
Einstein-$\Lambda$ tensor mode
$h_{k}=A\,(1-\ii k\eta)\,\ee^{\ii k\eta}
+B\,(1+\ii k\eta)\,\ee^{-\ii k\eta}$, their Eq.~(67).
In the proxy frame the translation is immediate:
\begin{itemize}
\item Eq.~(30), $\dot x=0$: the variable $q_{2}$, i.e.\ exactly
Maldacena's effect $E_{\sigma}$ of \eqref{eq:Esigma};
\item Eqs.~(59)/(64), $\ddot x=-\omega_{1}^{2}x$ (healthy Einstein
dispersion): the phase-space variable
\begin{equation}
O_{2}\;=\;\ddot x+\omega_{1}^{2}x\;=\;p_{2}+\omega_{1}^{2}q_{1},
\qquad\text{constrained to }0,
\label{eq:O2}
\end{equation}
regularized as $E'_{\tau}=\exp(-O_{2}^{2}/2\tau^{2})$. We adopt the
proxy dispersion $\ddot x=-\omega_{1}^{2}x$, which is the
\emph{ghost-killing} sign derived just below---not a free choice.
HLZ's Fourier-space condition (their Eq.~(64),
$\ddot h_{k}|_{0}=+k^{2}h_{k}|_{0}$, confirmed against the published
text) carries the opposite \emph{literal} sign yet selects the same
object: together with Neumann (their Eq.~(30), $\dot h|_{0}=0$) it
uniquely carves out the Einstein two-parameter family
$h_{k}=A(1-\ii k\eta)\ee^{\ii k\eta}+B(1+\ii k\eta)\ee^{-\ii k\eta}$
(their Eq.~(67)). At the future boundary this family has \emph{positive
curvature}, $\ddot h(0)=+k^{2}h(0)$ (locally
$h\simeq1+\tfrac12 k^{2}\eta^{2}$)---the boundary signature the second
condition enforces---while the corresponding healthy mode of the
\emph{flat} proxy is the oscillator $\ee^{\pm\ii\omega_{1}t}$, whose
curvature is $\ddot x=-\omega_{1}^{2}x$. It is this boundary-curvature
sign that inverts between the conformal-time and flat-proxy
descriptions---a property of the declared proxy, not a frame
convention---so the literal $+k^{2}$ becomes $-\omega_{1}^{2}$; the
proxy sign is fixed by ghost-killing (below). In
normal-mode language $O_{2}=0$ is precisely ``kill the ghost
coordinate'': for $x=y_{1}+y_{2}$ with
$\ddot y_{i}=-\omega_{i}^{2}y_{i}$, the condition
$\ddot x=-\omega_{1}^{2}x$ forces
$(\omega_{1}^{2}-\omega_{2}^{2})\,y_{2}=0$.
\end{itemize}
The crucial structural fact is that $p_{2}$ is the Ostrogradski
momentum conjugate to $q_{2}$:
\begin{equation}
[q_{2},O_{2}]=[q_{2},p_{2}]=\ii .
\label{eq:conjugate}
\end{equation}
Maldacena's $\mathcal{M}$ fixes one quadrature of the velocity sector; the HLZ
$\mathcal{M}'$ fixes a \emph{canonically conjugate pair}. $\mathcal{M}'$ is therefore a
joint position--momentum effect, and Heisenberg limits it: neither natural
ordering of the two Gaussian conditions can beat the minimal
phase-space cell. We price $\mathcal{M}'$ against the same reference vacuum
$\Psi_{0}$ as $\mathcal{M}$, so the two effects are scored on identical input;
HLZ's dropping of the in-state restriction constrains the
\emph{delivered} state, not the reference against which survival is
measured.

\begin{proposition}[Second-condition statistics]
\label{prop:O2}
On the vacuum of Proposition~\ref{prop:vacuum},
\begin{equation}
O_{2}\Psi_{0}=\bigl(\ii\,\delta\,q_{2}
+\omega_{1}\delta\,q_{1}\bigr)\Psi_{0},
\qquad
\mathrm{Var}_{\Psi_{0}}(O_{2})
=\frac{\delta}{2}\,\frac{\omega_{1}+\omega_{2}}{\omega_{2}}
\;\xrightarrow{\ \delta\to0\ }\;\delta,
\label{eq:varO2}
\end{equation}
and the symmetrized covariance vanishes exactly:
$\mathrm{Re}\,\langle q_{2}O_{2}\rangle_{\Psi_{0}}=0$
(indeed $\langle q_{2}O_{2}\rangle_{\Psi_{0}}=\ii/2$, pure
commutator). Consequently the $(q_{2},O_{2})$ Wigner marginal of $\Psi_{0}$ is an
uncorrelated Gaussian, so the \emph{Weyl}-ordered survival factorizes
as $p_{1}p_{2}$, with
\begin{equation}
p_{2}(\tau,\delta)
=\Bigl[1+\frac{\mathrm{Var}(O_{2})}{\tau^{2}}\Bigr]^{-1/2}
\;\xrightarrow{\ \delta\to0\ }\;1 .
\label{eq:p2}
\end{equation}
Operator orderings that respect $[q_{2},O_{2}]=\ii$ (product/Kraus)
do \emph{not} factorize; that non-factorization is precisely the
$1/(4\sigma^{2}\tau^{2})$ Heisenberg surcharge of
Proposition~\ref{prop:mprime}.
\end{proposition}

\begin{proof}
$p_{2}\Psi_{0}=-\ii\partial_{q_{2}}\Psi_{0}
=(\ii b q_{2}-cq_{1})\Psi_{0}$ and
$\omega_{1}^{2}-c=\omega_{1}(\omega_{1}-\omega_{2})
=\omega_{1}\delta$, giving the first identity. Then
$\mathrm{Var}(O_{2})=b^{2}\langle q_{2}^{2}\rangle
+\omega_{1}^{2}\delta^{2}\langle q_{1}^{2}\rangle
=\tfrac{b}{2}+\tfrac{\omega_{1}^{2}\delta^{2}}{2\alpha}
=\tfrac{\delta}{2}\bigl(1+\tfrac{\omega_{1}}{\omega_{2}}\bigr)$
(the cross terms cancel between the $\pm\ii$ pieces), and
$\langle q_{2}O_{2}\rangle
=\ii b\langle q_{2}^{2}\rangle=\ii/2$. Since $\Psi_{0}$ is Gaussian,
$O_{2}$ is Gaussian-distributed and \eqref{eq:p2} is the standard
Gaussian-effect integral. Grid values (spectral $p_{2}$,
$2048$-point FFT) reproduce \eqref{eq:varO2} to five digits at
$\delta=0.5,0.1,0.02$ and the covariance to $10^{-18}$ at
$\delta=0.3,0.05$.
\end{proof}

Equation \eqref{eq:varO2} says the prepared vacuum \emph{already
nearly satisfies} the second HLZ condition as $\delta\to0$: the
second condition is asymptotically free in scale. All the
$\delta$-collapse continues to come from the Neumann factor.
This conclusion---and with it the equal-leading-power verdict of
Proposition~\ref{prop:mprime}---rests on the sign
$O_{2}=p_{2}+\omega_{1}^{2}q_{1}$, which is the ghost-killing sign of
\eqref{eq:O2} and \emph{also} the finite-variance one: the opposite
sign $O_{2}=p_{2}-\omega_{1}^{2}q_{1}$ isolates no proxy mode and gives
$\mathrm{Var}(O_{2})=\tfrac{\delta}{2}
+\omega_{1}(\omega_{1}+\omega_{2})^{2}/(2\omega_{2}\delta)\to\infty$
(leading power $1$, not $\tfrac12$). Ghost-killing and finite variance
therefore select the same sign, and HLZ's Eq.~(64)---%
$\ddot h_{k}|_{0}=+k^{2}h_{k}|_{0}$, confirmed against the published
text---is consistent with it under the role inversion of
\eqref{eq:O2}. The second-condition sign is thus derived; the only
residual modeling input is the declared proxy dictionary itself
(Section~\ref{sec:limitations}).

\subsection{The price of \texorpdfstring{$\mathcal{M}'$}{M'}}

Because the two constrained variables are conjugate, the survival of
$\mathcal{M}'$ depends on the operator ordering of the two Gaussian factors---an
$O(1)$, not power-law, effect.

\begin{proposition}[$\mathcal{M}'$ price sheet]
\label{prop:mprime}
Let $p_{1}=p(\sigma,\delta)$ of \eqref{eq:price} and
$p_{2}=p_{2}(\tau,\delta)$ of \eqref{eq:p2}. Then:
\begin{itemize}
\item[(i)] (Weyl ordering; exact) The operator with Weyl symbol
$\exp(-q_{2}^{2}/2\sigma^{2}-O_{2}^{2}/2\tau^{2})$ has vacuum
expectation exactly $p_{1}p_{2}$; it is a bona fide effect
(positivity $E\ge0$) iff $\sigma\tau\ge1/2$. Here $\sigma\tau\ge1/2$ is the
Gaussian positivity/uncertainty condition securing the lower bound
$E\ge0$; the upper bound $E\le1$ then follows because at and above
the minimal cell the operator is, up to a metaplectic
(Gaussian-unitary) conjugation, a single-mode Gaussian effect whose
Fock eigenvalues lie in $[0,1)$, reducing to $\mathbf{1}$ only as
$\sigma,\tau\to\infty$.
\item[(ii)] (Product/Kraus ordering; exact) The always-legitimate
effect $E=K^{\dagger}K$ with
$K=\ee^{-O_{2}^{2}/4\tau^{2}}\allowbreak\ee^{-q_{2}^{2}/4\sigma^{2}}$ has vacuum
survival given exactly by the closed form
\begin{equation}
p_{\mathcal{M}'}\;=\;p_{1}\Bigl[1+\frac{\mathrm{Var}(O_{2})}{\tau^{2}}
+\frac{1}{4\sigma^{2}\tau^{2}}\Bigr]^{-1/2},
\label{eq:kraus}
\end{equation}
whose sub-minimal-cell (sharp-effect) limit, taken along the joint
$\sigma,\tau\to0$ path (so that the Heisenberg term
$1/(4\sigma^{2}\tau^{2})$ dominates the bracket, i.e.\
$\sigma\tau\to0$), is exactly
$p_{\mathcal{M}'}\to2\sigma\tau\,p_{1}$: the Maldacena survival times the
phase-space cell $2\sigma\tau$ in minimal-uncertainty units. The
extra $1/(4\sigma^{2}\tau^{2})$ in the bracket---absent from the Weyl
product $p_{1}p_{2}$ of (i)---is the Heisenberg back-reaction of
squeezing $q_{2}$ onto its conjugate $p_{2}\subset O_{2}$, i.e.\ the
$O(1)$ ordering surcharge.
\item[(iii)] Both orderings have the same leading detuning power,
exactly $1/2$: $p_{\mathcal{M}'}\sim\sigma\sqrt{2\delta}$ as $\delta\to0$ up to
$O(1)$ prefactors, since $p_{1}\to\sigma\sqrt{2\delta}$ while the
bracket in \eqref{eq:kraus} tends to a $\delta$-independent constant
(recall $\mathrm{Var}(O_{2})\to0$). The $\delta\to0$ leading power of
the closed form \eqref{eq:kraus} is therefore $1/2$ for every
$(\sigma,\tau)$.
\end{itemize}
\end{proposition}

\begin{proof}
(i) The expectation of a Weyl-ordered operator is the phase-space
average of its symbol against the Wigner function; the Wigner
function of the Gaussian $\Psi_{0}$, marginalized to the
$(q_{2},O_{2})$ plane, is an uncorrelated Gaussian with the variances
of Proposition~\ref{prop:O2}, so the double Gaussian integral
factorizes into $p_{1}p_{2}$. Positivity of the Weyl quantization of
a Gaussian symbol holds iff the symbol's covariance admits the
minimal-cell bound, i.e.\ $\sigma\tau\ge1/2$. The two single-factor
limits are exact on the grid: at $\delta=0.1$,
$\tau=10^{3}$ returns $p_{1}$ and $\sigma=10^{3}$ returns $p_{2}$,
each to five digits.

(ii)--(iii) In the product ordering
$\langle\Psi_{0}|K^{\dagger}K|\Psi_{0}\rangle
=\langle\phi|\ee^{-O_{2}^{2}/2\tau^{2}}|\phi\rangle$ with
$\phi=\ee^{-q_{2}^{2}/4\sigma^{2}}\Psi_{0}$, still Gaussian. Three
steps. (1)~The norm ratio is Maldacena's survival exactly,
$\norm{\phi}^{2}/\norm{\Psi_{0}}^{2}
=[1+1/(2\sigma^{2}\delta)]^{-1/2}=p_{1}$, since
$|\phi|^{2}=\ee^{-q_{2}^{2}/2\sigma^{2}}|\Psi_{0}|^{2}$ reproduces the
integral of Proposition~\ref{prop:survival}. (2)~In the normalized
state $\phi/\norm{\phi}$, $O_{2}$ is again a zero-mean Gaussian, now
with variance
$V'=\mathrm{Var}_{\Psi_{0}}(O_{2})+1/(4\sigma^{2})$: narrowing
$q_{2}$ by the factor $\ee^{-q_{2}^{2}/4\sigma^{2}}$ broadens the
conjugate $p_{2}\subset O_{2}$ by exactly $1/(4\sigma^{2})$, the
Heisenberg back-reaction of squeezing $q_{2}$. (3)~The
Gaussian-effect overlap then gives
$\langle\phi|\ee^{-O_{2}^{2}/2\tau^{2}}|\phi\rangle/\norm{\phi}^{2}
=[1+V'/\tau^{2}]^{-1/2}$. Multiplying (1) and (3) yields
\eqref{eq:kraus} exactly. Along the joint $\sigma,\tau\to0$ path (so
that $\sigma^{2}\mathrm{Var}(O_{2})\to0$) the
$1/(4\sigma^{2}\tau^{2})$ term dominates the bracket, giving
$p_{\mathcal{M}'}\to2\sigma\tau\,p_{1}$; and since $\mathrm{Var}(O_{2})\to0$
with $p_{1}\to\sigma\sqrt{2\delta}$, the detuning exponent is exactly
$1/2$. A high-resolution grid tracks \eqref{eq:kraus} (with
grid-resolution artifacts at the smallest $\delta$, where the vacuum
width $\sim\delta^{-1/2}$ outruns the box; the tabulated values are the
analytic closed form); the comparison table follows.
\end{proof}

\begin{table}[ht]
\centering
\begin{tabular}{cccccc}
\toprule
$\delta$ & $p_{1}$ & $p_{2}$ & $p_{1}p_{2}$ (Weyl) &
$p_{\mathcal{M}'}$ \eqref{eq:kraus} & $2\sigma\tau\,p_{1}$ \\
\midrule
$0.500$ & $0.09950$ & $0.11471$ & $0.01141$ & $0.00196$ & $0.00199$\\
$0.100$ & $0.04468$ & $0.29417$ & $0.01314$ & $0.00089$ & $0.00089$\\
$0.020$ & $0.02000$ & $0.57540$ & $0.01151$ & $0.00040$ & $0.00040$\\
$0.004$ & $0.00894$ & $0.84491$ & $0.00756$ & $0.00018$ & $0.00018$\\
\bottomrule
\end{tabular}
\caption{Pricing $\mathcal{M}'$ at $\sigma=\tau=0.1$
(so $2\sigma\tau=0.02$). The Neumann factor $p_{1}$ carries all the
$\delta$-collapse; the second-condition factor $p_{2}\to1$
(asymptotically free); the product-ordered survival is the exact
closed form \eqref{eq:kraus}, which in this deep sub-cell
(sharp-effect) regime tracks $2\sigma\tau\,p_{1}$, i.e.\ Maldacena's price times a
Heisenberg-cell surcharge. Here $\sigma\tau=0.01<1/2$, so by
Proposition~\ref{prop:mprime}(i) the Weyl operator is not positive in
this regime; the displayed $p_{1}p_{2}$ (Weyl) column is therefore the
expectation of a non-effect---a formal comparator, not a survival
probability.}
\label{tab:mprime}
\end{table}

\subsection{What the surcharge buys}

The fidelity scorecard exposes the qualitative difference between the
two exorcisms. On the same grids, at $\sigma=\tau=0.05$:
\begin{center}
\begin{tabular}{ccc}
\toprule
$\delta$ & $F_{\mathcal{M}}$ (Neumann only) & $F_{\mathcal{M}'}$ (both conditions) \\
\midrule
$0.500$ & $0.7996$ & $0.0635$ \\
$0.100$ & $0.5494$ & $0.0506$ \\
$0.020$ & $0.2744$ & $0.0508$ \\
$0.004$ & $0.1258$ & $0.0786$ \\
\bottomrule
\end{tabular}
\end{center}
$F_{\mathcal{M}}$ tracks the closed form \eqref{eq:fidelity}; $F_{\mathcal{M}'}$ is
roughly flat and small, and this is \emph{by design}, not by failure:
$\mathcal{M}'$ does not target the pure $\omega_{1}$ vacuum at all. It targets
the general two-parameter ($A,B$) Einstein solution with no in-state
preparation---a mixed delivered state whose overlap with any single
pure vacuum is correspondingly reduced (small in the numerics above,
not by a general bound). Pure-vacuum fidelity is the right
scorecard for $\mathcal{M}$ and the wrong one for $\mathcal{M}'$.

The comparison is then clean:
\begin{itemize}
\item \emph{Survival:} $p_{\mathcal{M}}\sim\sigma\sqrt{2\delta}$ and
$p_{\mathcal{M}'}\sim\sigma\sqrt{2\delta}$ up to $O(1)$: the second HLZ
condition adds only a Heisenberg surcharge (it fixes $\ddot x$,
conjugate to the $\dot x$ that Neumann fixes) and is not a new source
of collapse, because $\mathrm{Var}(O_{2})\sim\delta$ already
vanishes.
\item \emph{Delivery:} $\mathcal{M}$ delivers (at $\sqrt{\delta}$-collapsing
fidelity) a pure but over-broad vacuum; $\mathcal{M}'$ delivers the general
Einstein solution.
\end{itemize}
This is HLZ's own claim---Neumann alone yields no general solution---%
re-expressed as a price: the extra condition buys \emph{generality},
not economy, and the conformal ($\delta\to0$) collapse of the
survival is not evaded by either scheme. In the currency of this
series, there is no cheap exorcism of the conformal graviton pair by
future boundary data: the price is set by the exceptional-point
degeneration of the state, which no rearrangement of the boundary
conditions removes.

%=====================================================================
\section{The pole sheet is readable in the divisor}
\label{sec:sheet}

We now turn from boundary-condition exorcisms to
propagator-prescription exorcisms, and to the spectral question
underneath all of them: what, exactly, distinguishes a ghost from an
ordinary unstable particle with the same complex pole pair? Buoninfante
\cite{BuoninfanteP1,BuoninfanteP3} gives the sharp QFT answer: only
the \emph{sheet}. Both have poles at $m^{2}\pm\ii m\Gamma$; the
unstable particle's pair sits on the second (unphysical) sheet and its
K\"all\'en--Lehmann representation is a positive density on the real
axis with sum rule $1=C$ (complete norm transfer to the continuum:
the particle decays) \cite[Eq.~(34)]{BuoninfanteP3}; the ghost's pair
sits on the first (physical) sheet and the spectral decomposition
contains, in addition to a \emph{still-positive} continuum, discrete
negative-residue complex-pole terms
$-Z\Delta_{F}(m^{2}+\ii m\Gamma)-Z^{*}\Delta_{F}(m^{2}-\ii m\Gamma)$
\cite[Eq.~(40)]{BuoninfanteP3}, with sum rule $Z+Z^{*}=1+C$: the
one-particle component persists \cite[Eq.~(44)]{BuoninfanteP3}. We
emphasize the precise locus of positivity failure, because it is easy
to state wrongly: the ghost's continuum is positive; what fails is
positivity of the \emph{total} spectral measure, through the discrete
complex-pole part. The two survival regimes arise from two distinct
mechanisms, and it is worth keeping them separate. A positive
spectral measure on the real axis gives, by Bochner's theorem, a
\emph{bounded} amplitude $|A(t)|\le1$---the second-sheet resonance.
What makes the first-sheet ghost's $|A(t)|$ \emph{unbounded} is not
sign-indefiniteness as such---a finite signed measure on the real
axis still gives a bounded $|A(t)|$---but that its spectral support
includes \emph{complex} energies $m\pm\ii\gamma$ off the real axis,
whose $\ee^{\pm\gamma t}$ legs make $|A(t)|$ grow, as the divisor will
now report.

\subsection{The Bochner discriminant}

\begin{proposition}[Pole sheet $\to$ process class]
\label{prop:bochner}
Let $A(t)$ be the survival amplitude of a single mode.
\begin{itemize}
\item[(i)] (Second sheet; rigorous.) If
$A(t)=\int\ee^{-\ii Et}\,\dd\mu(E)$ for a positive probability
measure $\mu$ on $\mathbb{R}$---the case of an ordinary unstable
particle, whose dressed propagator has a positive spectral
density---then $A$ is a positive-definite function and Bochner's
theorem gives $|A(t)|\le A(0)=1$ for all $t$. The two-state survival
matrix
\[
\Gamma(t)=\begin{pmatrix}|A(t)|^{2}&0\\ 1-|A(t)|^{2}&1\end{pmatrix}
\]
is then column-stochastic with \emph{no conditioning}: because
$|A|\le1$, an unconditioned sub-stochastic survival process exists on the
second sheet whatever the negativity pattern. Divisor negativity over
$[t_{1},t_{2}]$ occurs iff $|A(t_{2})|>|A(t_{1})|$ (net increase over the
interval; backflow); its
magnitude $M_{21}=1-|A(t_{2})|^{2}/|A(t_{1})|^{2}$ is bounded only where
$\inf|A|>0$ and diverges near an amplitude null. For absolutely continuous
$\mu$ the Riemann--Lebesgue lemma sends $A(t)\to0$ in amplitude, so the
negativity is transient \emph{when $|A|$ is eventually monotone} (as for
the smooth density below)---the transient-indivisibility class of
\cite{indivisible}; a compactly supported or oscillatory a.c.\ measure
($\mu$ uniform on an interval, $A\propto\operatorname{sinc}$) instead
revives backflow at every amplitude null. A measure with atoms likewise
stays unconditioned ($|A|\le1$) without being transient---a single atom
over an a.c.\ background can even make $|A|$ eventually monotone
\emph{increasing} ($\mu=0.7\,\delta_{0}+\rho$ gives
$A=0.7-0.2\ee^{-t}+0.5\ee^{-2t}$, permanent second-sheet negativity with
no conditioning). What is invariant is the dichotomy by \emph{boundedness
of $|A|$}, not the negativity pattern: $|A|\le1$ always admits an
unconditioned process, while the unbounded $|A|$ of (ii) forces
conditioning.
\item[(ii)] (First sheet; mode-level analytic, closed forms
numerically re-verified.) If the mode dynamics carries a first-sheet
complex-conjugate pair---eigenfrequencies $m\pm\ii\gamma$,
$\gamma>0$, the mode realization of Buoninfante's ghost
pair---then no positive spectral measure exists for $A$, and indeed
$|A(t)|$ grows like $\cosh(\gamma t)$ on the configuration basis (any
beable state overlapping the growing leg): since $|A(t)|>1$ at finite
$t$, \emph{no} unconditioned sub-stochastic survival process exists
at any coarse-graining. The survival-conditioned dilation
\eqref{eq:conditioned} is mandatory, and for the dimer realization
\eqref{eq:dimer} it lands in the terminal class of
\cite{indivisible}: permanent divisor negativity past
$t_{c}$ of \eqref{eq:tcfloor}, with the late-time conditioned
generator pinned at the floor $-(a+\wt)$.
\end{itemize}
\end{proposition}

\begin{proof}
(i) is Bochner's theorem plus inspection:
$|A(t)|=|\int\ee^{-\ii Et}\dd\mu|\le\int\dd\mu=1$; column sums of
$\Gamma(t)$ are $1$; the divisor \eqref{eq:divisor} of the survival
pair has entries $M_{11}=|A(t_{2})|^{2}/|A(t_{1})|^{2}$ and
$M_{21}=1-M_{11}$, negative iff $|A|$ increased; Riemann--Lebesgue sends
$A\to0$ in amplitude, so the revivals decay when $|A|$ is eventually
monotone (but not for a compactly supported or oscillatory a.c.\ measure,
whose $|A|$ nulls revive backflow). A quantitative illustration: for the positive
density $\rho(E)\propto\sqrt{E}\,(1+E)^{-1/2}
\bigl[(E-m_{0})^{2}+\Gamma_{0}^{2}/4\bigr]^{-1}$
($m_{0}=3$, $\Gamma_{0}=0.4$), the computed
$\max_{t}|A(t)|=1.000000$, and the residual backflow episodes are at
the $10^{-9}$ level (numerical noise), with no growing revival---this particular density is
essentially monotone; the general argument allows densities with
sharper thresholds to show visible but decaying revivals, though the
public script tests only the smooth density above.
(ii) The growth statement is elementary: with eigenfrequencies
$m\pm\ii\gamma$ the evolution has singular values
$\ee^{\pm\gamma t}$ up to bounded factors, so
$\norm{U(t)}\ge C\ee^{\gamma t}$; on the dimer \eqref{eq:dimer} this
is exact via \eqref{eq:dimerU}. The computed
$\max_{t\le6}\norm{\ee^{-\ii Ht}}=1098.9$ at
$(a,\kappa)=(1.5,1)$. The terminal-class closed forms
\eqref{eq:tcfloor} are quoted from \cite{indivisible} and re-verified
here: permanent-negativity onset $0.6400$ on a $0.005$ time grid
against the closed form $0.6487$, and late-time generator
off-diagonal $-2.6031$ against the floor $-(a+\wt)=-2.6180$ (both
discrepancies are finite-difference artifacts that shrink with the
step). We state part (ii) at mode level: the assertion that the
\emph{field-theoretic} ghost two-point function has exactly this
discrete negative-residue structure is Buoninfante's
\cite[Eqs.~(40),(44)]{BuoninfanteP3}, not re-derived here.
\end{proof}

\begin{remark}[The classifier crosses at the sheet boundary]
Sweeping a two-mode model with eigenvalues $3+\ii g_{\mathrm{eff}}$
and $3-0.4\ii$ across the boundary: for
$g_{\mathrm{eff}}=-0.4,-0.1,0.0$ the maximal evolution norm is
$1.000$ (unconditioned, resonance class); at
$g_{\mathrm{eff}}=+0.1$ it is $20.1$ and at $+0.4$ it is
$1.6\times10^{5}$ (conditioning mandatory; not terminal---this sweep is
decoupled). The
process-class classifier changes exactly at the sheet boundary
$g_{\mathrm{eff}}=0$---there is no intermediate regime.
\end{remark}

The slogan form of Proposition~\ref{prop:bochner}:
\begin{quote}
\emph{Unstable particles survive unconditioned; ghosts survive only
conditioned. The pole sheet is readable in the divisor by the \emph{boundedness} of
the amplitude: an unconditioned process on the second sheet
($|A|\le1$), and---for the dimer---terminal negativity with a generator
floor on the first ($|A|$ unbounded).}
\end{quote}
This gives the divisor currency an operational meaning directly in
QFT terms: handed only the survival record of a mode, one can read
off which sheet its pole pair inhabits---from whether $|A|$ stays
bounded (second sheet) or grows without bound (first). The reading is
a converse only within the two-class family above, and for states
overlapping the growing leg: a signed real-axis measure also keeps
$|A|$ bounded, and a first-sheet mode probed in a decaying-leg-only
state leaves a bounded record.

\subsection{The masking law is the terminal-class dimer's cross-population}
\label{sec:masking}

Buoninfante's second and third papers establish what he calls
\emph{anti-instability}: in the indefinite-metric theory the ghost
neither decays (norm conservation forbids negative-norm $\to$
positive-norm transfer) nor runs away; instead the would-be decay
manifests as \emph{multi-particle masking}, an overlap with
multi-particle states that grows as
$\sinh^{2}(\Gamma t/2)$: the general overlap is his
\cite[Eq.~(84)]{BuoninfanteP2}, whose at-rest ($\vec p=0$)
narrow-width form \cite[Eq.~(88)]{BuoninfanteP2} is exactly
$\sinh^{2}(\Gamma t/2)+O(\Gamma^{2})$, on the timescale $2/\Gamma$; no free
asymptotic one-particle ghost state exists, and the Feynman ghost
\emph{appears} to retain a single arrow of causality on shell---the
source's own hedge, which flags that first-sheet complex-conjugate
poles could still induce violations.

\begin{proposition}[Masking law $=$ cross-population growth]
\label{prop:masking}
For the broken dimer \eqref{eq:dimer}, the cross-population is
exactly
\[
|U_{12}(t)|^{2}=\frac{\kappa^{2}}{\wt^{2}}\,\sinh^{2}(\wt t),
\]
by \eqref{eq:dimerU}. Setting $\wt=\Gamma/2$ this is Buoninfante's
masking law $\sinh^{2}(\Gamma t/2)$ up to the constant prefactor
$(\kappa/\wt)^{2}$, including the timescale $2/\Gamma$. (The identity
is exact---algebraic, not fitted.)
\end{proposition}

The identification says the two formalisms are computing the same
growth and booking it differently. In Buoninfante's closed,
indefinite-metric currency, nothing decays and nothing runs away:
the growth is an overlap, and ``anti-instability'' is the statement
that the indefinite norm is conserved. In the beable currency of this
series, the same growth is the reason no unconditioned process
exists, and conditioning divides it out exactly: \emph{after}
conditioning, nothing physical grows---which is the conditioned
transcription of anti-instability. Two bookkeepings, one pole pair,
one shared closed form. We are careful about scope: Buoninfante's
masking is an overlap statement in indefinite-metric QFT; our reading
adds the sample-space price tag and claims nothing about his
finite-time resummation---what this paper adds is the currency, not a
new field theory.

%=====================================================================
\section{A taxonomy of exorcisms and their currencies}
\label{sec:taxonomy}

Buoninfante's first paper organizes the propagator prescriptions for
a ghost resonance into an $(a;b,c)$ taxonomy
\cite[Eqs.~(5)--(6)]{BuoninfanteP1}: $a=-1$ selects the ghost (first
sheet) and $a=+1$ the non-ghost/unstable particle (second sheet, cf.\
Section~\ref{sec:sheet}), and $(b,c)$ select how the complex pair is
routed through the contour. The three standard exorcisms, and the
unstable-particle control, then line up against the trilemma of
\cite{choosing}---a coupled-ghost theory keeps at most two of
\{coupled ghosts, empirical Lorentz invariance, definite
configurations\}---as follows.

\begin{table}[ht]
\centering
\footnotesize
\begin{tabular}{lllll}
\toprule
Prescription & Norms & Causal structure & Process reading &
Currency paid \\
\midrule
Feynman ghost & indefinite & causal &
conditioning mandatory, & sample space \\
($b=c=+1$) & (odd-ghost $<0$) & &
terminal class (\S\ref{sec:sheet}) &
(no beable seat) \\[3pt]
anti-Feynman & positive & acausal &
unconditioned, & causal order / \\
($b=c=-1$) & & (reversed arrow) &
reversed arrow & Lorentz \\[3pt]
fakeon & --- & purely virtual &
no asymptotic state & silence \\
($b=-c$) & & ($\mathrm{Im}=0$) & to condition &
(excision) \\[3pt]
unstable particle & positive & causal &
unconditioned & none \\
($a=+1$, 2nd sheet) & & & sub-stochastic & \\
\bottomrule
\end{tabular}
\caption{The $(a;b,c)$ prescription taxonomy of
\cite{BuoninfanteP1} mapped to the trilemma currencies of
\cite{choosing}. Every exorcism pays in at least one currency; only
the genuine second-sheet unstable particle is free.}
\label{tab:taxonomy}
\end{table}

Three remarks, in increasing order of contentiousness.

\emph{First}, the Feynman ghost row is
Proposition~\ref{prop:bochner}(ii) verbatim: keeping the ghost
coupled, causal, and on the first sheet forfeits the unconditioned
sample space; the theory survives in the beable currency only as a
survival-conditioned terminal-class process. This is the
``exit-the-sample-space'' horn of the trilemma, now anchored in the
sheet structure rather than asserted.

\emph{Second}, the anti-Feynman row is structurally the
Donoghue--Menezes route \cite{DonoghueMenezes}: a positive-norm
\emph{asymptotic spectrum} and a
unitary $S$-matrix are retained at the cost of microcausality (their
unstable-ghost loops arrow the ghost propagation backwards). We note
explicitly that this alignment is ours to draw: Buoninfante's trilogy
does not itself connect the $(b=c=-1)$ contour routing to the
``Merlin'' treatment of \cite{DonoghueMenezes}; the alignment is
between that routing and
the acausal-but-unitary corner they occupy,
and it slots into the Lorentz/causal-order column of the trilemma
table (the currency whose empirical price \cite{instantaneous}
computes for coupled-ghost vacuum decay).

\emph{Third}, the fakeon row is Anselmi's, and here sits the live
dispute. Anselmi, Briscese, Calcagni, and Modesto
\cite{Anselmi2025} argue that among the prescriptions for complex
poles only the fakeon is viable; Buoninfante's anti-instability
result argues the Feynman ghost is consistent as it stands
\cite{BuoninfanteP2,BuoninfanteP3}. Read in
Table~\ref{tab:taxonomy}, the two camps are not disagreeing about a
computation---they are pricing in different currencies. The fakeon
buys consistency with silence: the ghost never appears as an
asymptotic state, at the price of excising it from the sample space
altogether (there is nothing left to condition on). Buoninfante's
Feynman ghost buys consistency with conditioning: the ghost persists
(sum rule $Z+Z^{*}=1+C$), is masked rather than decayed, and---in our
transcription---its beable process is survival-conditioned, terminal
class. Each camp's ``viability'' is consistency in its own currency.
The trilemma does not adjudicate which currency is the right one to
pay; it prices both, and this paper's contribution is that the
\emph{boundary-condition} prices are now numbers
(Sections~\ref{sec:maldacena}--\ref{sec:sheet}) rather
than slogans.

%=====================================================================
\section{Relation to final-state post-selection}
\label{sec:priorart}

The general idea that a future boundary condition is a
post-selection is not ours. It is the engine of the
Horowitz--Maldacena black-hole final-state proposal
\cite{HorowitzMaldacena}, where a maximally entangled final state at
the singularity post-selects the interior dynamics; and it has been
proposed in cosmology, where Anastopoulos reads cosmic acceleration
itself as a quantum post-selection effect \cite{Anastopoulos2024}.
The two-state-vector and decoherent-histories traditions provide the
general formalism. What we have not found anywhere is the
\emph{intersection} with ghost removal: no reading of Maldacena's
Neumann condition \cite{Maldacena2011} (or the HLZ variant
\cite{HLZ}) as a survival-conditioned quantum effect, and no
computation of the survival/fidelity price of any ghost-removing
boundary condition. Our scope claims are therefore limited to
exactly that intersection: (i) the identification of ghost-removal
boundary conditions as effects in the technical sense, with the
state-prep/effect split of Section~\ref{sec:effect}; and (ii) the
quantitative price sheet, Tables~\ref{tab:pricesheet}
and~\ref{tab:mprime}, with its conformal collapse. Likewise, on the
prescription side, the pole-sheet/process-class dictionary of
Proposition~\ref{prop:bochner} and the masking identification of
Proposition~\ref{prop:masking} are, to our knowledge, new, while the
spectral inputs are entirely Buoninfante's
\cite{BuoninfanteP1,BuoninfanteP2,BuoninfanteP3}.

%=====================================================================
\section{Limitations}
\label{sec:limitations}

Every caveat below is substantive; none is decorative.

\begin{enumerate}
\item \emph{The sharp effect is improper.} The exact Neumann
projection has survival probability zero; all prices are stated for
the regularized Gaussian family \eqref{eq:Esigma} at fixed
$\sigma$, with the $\sigma\to0$ asymptotics displayed. The
$\delta$-scaling---the physical content---is uniform in $\sigma$.
\item \emph{The quantization is a declared choice.} We price the
Dirac-normalizable (beable-carrying) quantization of the PU mode
(Proposition~\ref{prop:vacuum}). The contour-rotated Bender--Mannheim
quantization \cite{BenderMannheim2008} is a different leg of the
taxonomy, and the price of the same boundary condition there is a
separate (open) computation; the answer may differ.
\item \emph{de~Sitter only; Minkowski is the singular limit.} The
proxy's clean graviton~$\otimes$~partner split exists only at
$\delta>0$; flat space is the Jordan-block point where the effect
$\mathcal{M}=\mathbf{1}\otimes E$ is not even well defined. We treat
$\delta\to0$ as the singular limit and verify that the price sheet
detects it (Corollary~\ref{cor:collapse})---that the currency
diagnoses the obstruction is a feature, but it also bounds the
method: nothing here prices a flat-space exorcism, because on this
analysis there is no well-posed one to price.
\item \emph{No double-counting, and the split is a modeling
decision.} We count positive frequency as preparation and Neumann as
the effect. This matches the logical roles in
\cite{Maldacena2011}, but a reader who insists on counting the
in-state restriction as part of the removal would obtain a different
(strictly smaller) survival; our convention is the charitable one.
\item \emph{Tree level only.} The whole analysis, like the trilemma
it extends \cite{choosing}, is tree-level. Beyond tree level a
future-boundary effect inherits the known interaction obstruction of
final-state proposals \cite{HorowitzMaldacena}---interactions between
the post-selected sector and spectators degrade the conditioning---%
and we make no claim there.
\item \emph{The proxy is declared, not derived.} The detuned PU mode
with $\delta$ standing for the $\Delta=3/\Delta=2$ splitting is a
mode-level translation of \cite{Maldacena2011}, adopted for its
exact solvability. Its qualitative structure (four solutions, the
$t$-growing partner, the Jordan degeneration at flat space) matches
the graviton system, and its prediction reproduces the HLZ
flat-space obstruction; but the numerical prefactors in
Tables~\ref{tab:pricesheet}--\ref{tab:mprime} are properties of the
proxy, not of conformal gravity. In particular, where
\cite{Maldacena2011}'s construction is an exact tree-level identity, the
proxy delivers a fidelity $F<1$ at $O(1)$ $\delta$ (e.g.\ $F\approx0.80$
at $\delta=0.5$): the proxy has no counterpart of the renormalization
step that makes that identity exact, so the $O(1)$ part of the deficit
is a proxy artifact. Moreover the $\delta\to0$ collapse of the
\emph{fidelity} specifically is reference-convention-dependent (a
co-scaling Einstein reference removes it; text after
Proposition~\ref{prop:delivery}); the convention-free claimed content
is the survival collapse $p\sim\sigma\sqrt{2\delta}$ and the delivered
state's absolute width divergence $\alpha^{-1/2}\sim\delta^{-1/2}$.
\item \emph{One result is a numerical re-verification of closed
forms.} The terminal-class landing of the conditioned first-sheet
process uses the critical time and generator floor proved in closed
form in \cite{indivisible}, re-verified here numerically; it is
presented at exactly that strength in
Proposition~\ref{prop:bochner}.
\end{enumerate}

%=====================================================================
\section{Conclusion}
\label{sec:conclusion}

Every exorcism has a price, and the divisor currency makes it
legible.

Maldacena's future boundary condition, the cleanest ghost removal in
the literature, is per mode a Gaussian survival effect on a beable
variable, and its full price sheet is closed form: survival
$\sigma\sqrt{2\delta}$, Einstein fidelity $2\sqrt{\omega_{2}\delta}$
(in units $\omega_{1}=1$, against a fixed Einstein reference),
joint figure of merit linear in the de~Sitter splitting. The
exorcism is affordable exactly where Maldacena performed it---in
de~Sitter, at $\delta=O(1)$---and its price diverges at the conformal
point, which is the flat-space obstruction of Hell, L\"ust, and
Zoupanos rediscovered as a collapse of the price sheet. Their own
two-condition variant fixes a canonically conjugate pair and is
therefore Heisenberg-limited: same leading price, an $O(1)$
surcharge, and what the surcharge buys is generality (the full
Einstein solution space) rather than economy. There is no cheap
version: the cost is set by the exceptional-point degeneration of the
state itself.

On the propagator side, the pole sheet is readable in the divisor.
Positive spectral measure---the second-sheet unstable particle---%
gives $|A(t)|\le1$ by Bochner's theorem and an unconditioned
sub-stochastic process whose divisor negativity decays when the
amplitude is eventually monotone decreasing (transient-indivisibility class). The first-sheet ghost pair, whose spectral decomposition
carries discrete negative-residue complex poles with the persistence
sum rule $Z+Z^{*}=1+C$, admits no unconditioned process at all:
conditioning is mandatory; for the dimer realization it is terminal, with the critical time and
generator floor of the first companion paper. Buoninfante's masking
law is the \emph{time dependence} of the terminal-class dimer's
cross-population---the
pre-conditioning growth that conditioning divides out---up to the
normalization $(\kappa/\wt)^{2}$; his
anti-instability and our conditioning are the closed-currency and
beable-currency descriptions of one growth. And the prescription
taxonomy---Feynman ghost, anti-Feynman, fakeon---pays out in the three
trilemma currencies: sample space, causal order, silence. The
Anselmi--Buoninfante dispute over which prescription is viable is a
dispute over which currency one is willing to pay in; both sides are
priced here.

The unifying statement is the one the title promises. A ghost that is
removed has not been annihilated; it has been \emph{paid for}---%
conditioned on survival of a boundary effect, or bought out by a
contour choice or an excision---and the price always shows up in the
books: as a vanishing survival probability, as a fidelity collapse,
as an abandoned causal order, or as a sample space with one fewer
inhabitant. Unstable particles survive unconditioned. Ghosts survive
only by paying---conditioned, or in another trilemma currency. For the
boundary-condition exorcisms priced here, that price is now a number.

%=====================================================================
\phantomsection
\section*{Author contributions and use of AI}
\addcontentsline{toc}{section}{Author contributions and use of AI}

This manuscript was produced through a collaboration, directed by the
author, with the large language model Claude Fable~5 (Anthropic),
which carried out the technical work: the research program, the
derivations and proofs, the numerical experiments, the literature
analysis, and the text, including this statement. J.T.\ posed the
research direction, provided iterative review, and takes
responsibility for the decision to circulate. The research notebook
in the program's public repository records the development of this
paper's results, including the scoping analyses in which the
state-preparation/effect split and the Heisenberg-conjugate structure
of the two-condition scheme were identified before the prices were
computed.

%=====================================================================
\phantomsection
\section*{Code availability}
\label{sec:code}
\addcontentsline{toc}{section}{Code availability}

All numerical claims are reproducible from three standalone scripts
(NumPy/SciPy, seconds to minutes each, controls printed inline) in
the paper's repository:
\begin{center}
\url{https://github.com/embeddetech/p08-exorcised-ghosts}
\end{center}
(also archived, with the full program history, in the hub repository
\url{https://github.com/embeddetech/ghosts-program}):
\begin{itemize}
\item \texttt{scripts/p8\_maldacena\_price.py}
(Section~\ref{sec:maldacena}): vacuum verification, survival and
fidelity closed forms, the price sheet;
\item \texttt{scripts/p8\_hlz\_mprime.py}
(Section~\ref{sec:hlz}): $\mathrm{Var}(O_{2})$, covariance,
ordering-dependent survival, delivered-state fidelities;
\item \texttt{scripts/p8\_buoninfante\_map.py}
(Section~\ref{sec:sheet}): the Bochner bound, the terminal-class
closed forms, the sheet-crossing sweep.
\end{itemize}

%=====================================================================
\begin{thebibliography}{99}\small

\bibitem{indivisible}
J.~Torkelson,
\emph{Indivisible stochastic processes for non-Hermitian and
higher-derivative quantum systems: detailed balance, broken $\PT$
symmetry, and the explosion of ghosts},
companion paper (2026),
\url{https://github.com/embeddetech/p01-indivisible-ghosts}.

\bibitem{instantaneous}
J.~Torkelson,
\emph{Ghost vacuum decay as a stochastic process: instantaneous
states, the boost-orbit measure, and the foliation that ghosts make
visible},
companion paper (2026),
\url{https://github.com/embeddetech/p03-instantaneous-ghosts}.

\bibitem{choosing}
J.~Torkelson,
\emph{Conformal gravity must choose: exceptional-point kinematics,
vacuum decay, and the sample space of Weyl gravity},
companion paper (2026),
\url{https://github.com/embeddetech/p05-choosing-ghosts}.

\bibitem{Maldacena2011}
J.~Maldacena,
\emph{Einstein gravity from conformal gravity},
arXiv:1105.5632 (2011).

\bibitem{HLZ}
A.~Hell, D.~L\"ust, and G.~Zoupanos,
\emph{On the ghost problem of conformal gravity},
JHEP \textbf{08} (2023) 168, arXiv:2306.13714.

\bibitem{BuoninfanteP1}
L.~Buoninfante,
\emph{Remarks on ghost resonances},
JHEP \textbf{02} (2025) 186, arXiv:2501.04097.

\bibitem{BuoninfanteP2}
L.~Buoninfante,
\emph{Asymptotic quantum dynamics of ghost fields},
arXiv:2605.29047 (2026).

\bibitem{BuoninfanteP3}
L.~Buoninfante,
\emph{Ghosts versus unstable particles in quantum field theory},
arXiv:2606.18349 (2026).

\bibitem{Anselmi2025}
D.~Anselmi, F.~Briscese, G.~Calcagni, and L.~Modesto,
\emph{Amplitude prescriptions in field theories with complex poles},
JHEP \textbf{05} (2025) 145, arXiv:2503.01841.

\bibitem{DonoghueMenezes}
J.~F.~Donoghue and G.~Menezes,
\emph{Unitarity, stability, and loops of unstable ghosts},
Phys.\ Rev.\ D \textbf{100}, 105006 (2019), arXiv:1908.02416.

\bibitem{HorowitzMaldacena}
G.~T.~Horowitz and J.~Maldacena,
\emph{The black hole final state},
JHEP \textbf{02} (2004) 008, arXiv:hep-th/0310281.

\bibitem{Anastopoulos2024}
C.~Anastopoulos,
\emph{Final states in quantum cosmology: cosmic acceleration as a
quantum postselection effect},
Phys.\ Rev.\ D \textbf{111}, 064045 (2025), arXiv:2401.07662.

\bibitem{PaisUhlenbeck1950}
A.~Pais and G.~E.~Uhlenbeck,
\emph{On field theories with non-localized action},
Phys.\ Rev.\ \textbf{79}, 145 (1950).

\bibitem{BenderMannheim2008}
C.~M.~Bender and P.~D.~Mannheim,
\emph{No-ghost theorem for the fourth-order derivative
Pais--Uhlenbeck oscillator model},
Phys.\ Rev.\ Lett.\ \textbf{100}, 110402 (2008), arXiv:0706.0207.

\end{thebibliography}

\end{document}
